%vsgtu709



 %%%

%%%   Самарский государственный технический университет

%%%   Кафедра прикладной математики и информатики

%%%

%%%   Преамбула для статьи в журнал "Вестник Самарского государственного

%%%   технического университета. Серия: «Физико-математические науки»"

%%%   Ориентирована под MikTEx 2.4 for Win32

%%%

%%%   Хотя сам журнал собирается под GNU/Linux на дистрибутиве TeXLive



\documentclass[11pt,a4paper,twoside]{article}



\usepackage[cp1251]{inputenc}

\usepackage[T2A]{fontenc}

\usepackage[english,russian]{babel}

\usepackage{samgtubib}

\usepackage[dvips]{graphicx}

\usepackage[multiple]{footmisc}



\usepackage{samgtu}

\renewcommand{\VestnikYear}{2009} % Год издания Вестника

\renewcommand{\VestnikNumber}{2\,(19)} % Номер Вестника

\renewcommand{\VestnikMonth}{Ноябрь} % Месяц выхода Вестника

\renewcommand{\VestnikData}{01.11.09} % Дата подписания Вестника в печать

\renewcommand{\VestnikUPL}{26,64} % Усл. печ. л.

\renewcommand{\VestnikUIL}{25,73} % Уч.-изд. л.

\renewcommand{\VestnikRegNumber}{479/09} % Регистрационный номер

\renewcommand{\VestnikOrder}{994} % Номер заказа







%

 \begin{document}

 \tableofcontents

 \newpage



\renewcommand{\firstpage}{0}

\renewcommand{\lastpage}{0}

 \newcommand{\bigcom}{\text{\usefont{T2A}{cmfib}{m}{n}\large{\raisebox{1pt}{,}}}} % Cимвол производной в ТФКП (ковариантная)

%\Partition{Математическое моделирование, численные методы и комплексы программ}{Applied mathematics} % Раздел Вестника, в который подаётся статья

% Названия разделов см. на сайте при подаче заявки





%%%%%%%%%%%%%%%%%%%%%%%%%%%%%%%%%%%%%%%%%%%%%%%%%%%%%%%%%%%%%%%%%%%%%%%%%%%%%%%%%%%%

%Информация о статье и об авторах на русском и английском языках



% Номер УДК

\udk{532.517.2, 532.546}%Обработка текста. Подготовка текстов : Языки программирования. Компьютерные языки

\msc{76D33}% Коды MSC см. http://www.ams.org/msc/



\title{Квазиодномерная модель гиперболического типа развития трещины гидроразрыва пласта}% Полный заголовок

{Квазиодномерная модель гиперболического типа развития трещины гидроразрыва пласта \dots}% Заголовок, который идёт в верхний колонтитул



\author{А.\,М.~Ильясов, Г.\,Т.~Булгакова}% Автор(ы) для заголовка, если авторов несколько укать \inst

{И\,л\,ь\,я\,с\,о\,в~А.\,М.,\ Б\,у\,л\,г\,а\,к\,о\,в\,а~Г.\,Т.}% Автор(ы) для оглавления и колонтитула через \,



\institute{ФГБОУ ВО «Уфимский государственный авиационный технический университет», 450008, Уфа, К. Маркса, 12.}

% Если несколько, то так  \institute{ Организация 1 \and Организация 2  \and Организация 3}



\email{E-mails}{amilyasov67@gmail.com, bulgakova.guzel@mail.ru} %E-mail's авторов



\otherinf{\fullauthor{Ильясов Айдар Мартисович}\regalia{Уфа, Россия }\position{кандидат физико-математических наук, старший научный сотрудник} \dept{кафедры математики УГАТУ (450008, г. Уфа, К. Маркса, 12, e-mail: amilyasov67@gmail.com)}

\fullauthor{\\Булгакова Гузель Талгатовна}\regalia{Уфа, Россия }\position{доктор физико-математических наук, профессор, профессор} \dept{кафедры математики УГАТУ (450008, г. Уфа, К. Маркса, 12, e-mail: bulgakova.guzel@mail.ru, автор, ведущий переписку)} }



\keywords{гидроразрыв пласта, характеристики, инварианты Римана, эволюционность границ трещины}



\workabstract{% Здесь должна быть краткая аннотация статьи, например такая:

Представлена квазиодномерная модель гиперболического типа развития трещины гидроразрыва пласта (ГРП) в предположении, что коэффициент интенсивности напряжений намного превосходит коэффициент трещиностойкости породы. Рассматриваемая модель, учитывающая конвективные и нестационарные слагаемые в уравнении движения жидкости, является обобщением локальной модели Перкинса – Керна – Нордгрена (модель PKN). Доказано, что полученная система дифференциальных уравнений является квазилинейной строго гиперболической системой, для которой найдены характеристики и соотношения на характеристиках. В случае пренебрежения поправкой Кориолиса найдены инварианты Римана. При пренебрежении фильтрационной утечкой и вязкостью закачиваемой жидкости определены волны Римана, аналогичные простым плоским волнам в газовой динамике и изучены их свойства. Исследована эволюционность границ трещины ГРП. Поставлена начально-краевая задача развития трещины ГРП. Показано, что при пренебрежении диссипативными слагаемыми в представленной модели можно построить теорию простых волн, аналогичную теории одномерной газовой динамики с изоэнтропическими плоскими волнами.}



\engtitle{The quasi-one-dimensional hyperbolic model of hydraulic fracturing}



\engauthor{A.\,M.~Ilyasov, G.\,T.~Bulgakova}{I\,l\,y\,a\,s\,o\,v~A.\,M.,

B\,u\,l\,g\,a\,k\,o\,v\,a~G.\,T.}





\enginstitute{Ufa State Aviation Technical University,

12, K.Marx st., Ufa, 450008, Russian Federation

}% Место работы каждого из авторов на англ. языке

% Если несколько, то так  \enginstitute{ Организация 1 \and Организация 2  \and Организация 3}





\otherenginfm{\fullauthor{Aidar M. Ilyasov } \regalia{Ph.\,D. (Phys. \& Math.)} \position{Ass. Professor}

\dept{Department of Mathematics, Ufa State Aviation Technical University (450008, 12, K. Marx st., Ufa, Russian Federation, e-mail: bulgakova. guzel@mail.ru) }

\fullauthor{\\Guzel T. Bulgakova } \regalia{Ph.\,D. (Phys. \& Math.)} \position{Professor}

\dept{Department of Mathematics, Ufa State Aviation Technical University (450008, 12, K. Marx st., Ufa, Russian Federation, e-mail: amilyasov67@gmail.com)}}



\engabstract{The paper describes a quasi-one-dimensional hyperbolic model of hydraulic fracture growth assuming for the hydraulic fracturing that stress intensity is much higher than fracture resistance. The mode under analysis, which accounts for convective and unsteady terms in the fluid flow equation, is a generalization of the Perkins - Kern - Nordgren local model (PKN model). It has been proved that the obtained system of differential equations is a quasi-linear strictly hyperbolic system, for which the characteristics were found as well as their correlations. For the case of the Coriolis correction neglect, the Riemann invariants were found. Neglecting the injected fluid leak-off and viscosity, the Riemann waves, similar to simple plane waves in gas dynamics, were defined and their properties were studied. The evolutionism of fracture boundaries was investigated. The initial boundary value problem was set for fracture growth. It has been shown that the neglect of dissipative terms in the presented model allows constructing a simple wave theory analogous to the theory of one-dimensional gas dynamics for isentropic plane waves.  }



\engkeywords{hydraulic fracturing, characteristics, Riemann invariants, the evolutionism of fracture boundaries.\\ Булгакова Гузель Талгатовна: http://orcid.org/0000-0001-8030-1791. \\Ильясов Айдар Мартисович:   http://orcid.org/0000-0002-4399-9040

}





%%%%%%%%%%%%%%%%%%%%%%%%%%%%%%%%%%%%%%%%%%%%%%%%%%%%%%%%%%%%%%%%%%%%%%%%%%%%%%%%%%%%





\maketitle



%Каждый пункт статьи должен начинаться с команды Название пункта

%после названия точку ставить не нужно. Пункты автоматически нумеруются.

%Если пункт нумеровать не нужно, то следует использовать в команде

% \Section необязательный параметр [n]:

%   \Section[n]{Имя пункта}

% Введение и заключение обычно не нумеруются.



\Section[n]{Введение}

В настоящей работе моделируется режим развития трещины ГРП, в котором коэффициент интенсивности напряжений намного превосходит коэффициент трещиностойкости породы. В результате происходит разрушение породы на масштабах порядка размера пор в виде трещины-предвестника, фронт которой намного опережает фронт трещины ГРП, которая образуется вследствие нагнетания жидкости разрыва. В этом случае раскрытие вертикальной трещины ГРП происходит вследствие превышения давлением жидкости минимального горизонтального напряжения в породе и последующего заполнения уже существующей трещины-предвестника, образовавшейся в результате разрушения породы. Основы теории развития вертикальных трещин в пористой среде заложены в работе$ ~\cite{Zheltov-Christ} $.



В последние два десятилетия, как в отечественной, так и зарубежной литературе появилось более двух десятков тысяч теоретических и прикладных работ по моделированию образования и развития трещин ГРП. В данном обзоре рассматриваются только квазиодномерные модели развития трещины ГРП, применение которых широко распространено в нефтедобывающей промышленности. В этих моделях жидкости гидроразрыва предполагаются ньютоновскими и несжимаемыми, а окружающий пористый массив считается линейно-упругим. Вследствие этого в одномерных моделях возникает линейная связь между давлением в трещине и ее раскрытием. По характеру этой связи различают локальные и нелокальные модели. В локальных моделях между раскрытием и давлением имеется конечная связь, а в нелокальных моделях эта связь представляет собой сингулярное интегральное уравнение$~\cite{Mush}$. Теоретические исследования этих моделей в основном посвящены нахождению автомодельных решений. В работах$~\cite{Perkins_Kern, Nordgren1, Iv_Sm1, Basn1}$  изучаются локальные модели, а в работах$~\cite{Teod_tr1,Garagash1,Garagash2, Adachi1,Mitchell1}$ – нелокальные модели.



Модель распространения вертикальных трещин в пористой среде в случае ньютоновской жидкости гидроразрыва, когда длина трещины много больше ее постоянной высоты и ширины разработана в работах$~\cite{Perkins_Kern, Nordgren1}$ – локальная модель Перкинса – Керна – Нордгрена (модель PKN). В модели PKN течение в пласте не рассматривается и вводится эмпирический коэффициент фильтрационной утечки жидкости в пласт через стенки трещины. В рамках этой модели в работе$~\cite{Nordgren1}$ получены аналитические автомодельные решения в случаях малой и большой утечек, а также численные решения в общем случае.



В работе$~\cite{Iv_Sm1}$ рассмотрено обобщение классической модели PKN с зоной пропитки с прямолинейным фильтрационным течением, зависящим только от времени в несжимаемом пласте. Получены автомодельные решения в предельных случаях малых и больших фильтрационных утечек при заданном постоянном расходе или постоянном давлении на входе в трещину. В этих предельных случаях степенные зависимости искомых функций от времени близки к аналогичным зависимостям, полученным в$~\cite{Nordgren1}$. Также в$~\cite{Iv_Sm1}$ показано, что в общем случае, автомодельное решение возможно только при линейной зависимости давления от времени на входе в трещину.



В работе$~\cite{Gord_Z1}$ предложена двухзональная модель PKN, которая обобщает модель, представленную в$~\cite{Iv_Sm1}$. В работе$~\cite{Gord_Z1}$ кроме зоны пропитки также учитывается изменение порового давления в дальней зоне пласта, где предполагается упругий режим фильтрации пластовой жидкости~\cite{Smirn_Tag1}. Методом последовательных приближений с применением многочленов Эрмита получено автомодельное решение при постоянном давлении на входе в трещину.



В работе~\cite{Basn1} рассмотрена та же модель PKN с зоной пропитки, что и в~\cite{Iv_Sm1}. Получены автомодельные решения степенного и экспоненциального вида в предельных случаях малых и больших фильтрационных утечек при расходе или давлении на входе в трещину такого же вида. В~\cite{Basn1} показано, что в общем случае конечной фильтрационной утечки, автомодельное степенное решение возможно только при линейной зависимости давления от времени на входе в трещину или квадратичной зависимости расхода от времени.



В работе~\cite{Teod_tr1} рассмотрена модель PKN с нелокальной связью раскрытия и давления в трещине. Получено автомодельное решение для раскрытия трещины и давления в виде разложения по полиномам Чебышева второго рода при постоянном расходе на входе и отсутствии фильтрационной утечки. Показано, что при умеренных расходах закачки форма трещины близка к эллиптической форме. При больших расходах градиент давления на конце трещины увеличивается, а длина трещины – уменьшается.



В рамках нелокальной модели с геометрией трещины, аналогичной геометрии классической модели PKN, методами асимптотических разложений получены полуаналитические автомодельные решения для различных предельных режимов течения в трещине. В работе~\cite{Garagash1} получено полуаналитическое решение для режима течения в трещине с непроницаемыми стенками в пренебрежении коэффициентом трещиностойкости породы. В работе~\cite{Garagash2} получено решение для режима течения в трещине с непроницаемыми стенками с высоким коэффициентом трещиностойкости породы. В отличие от работ~\cite{Garagash1, Garagash2} в работе~\cite{Adachi1} представлен численный алгоритм для режима течения с фильтрационной утечкой, но с малой трещиностойкостью породы. Было проведено сравнение полученного численного решения на малых и больших временах с соответствующими асимптотическими решениями. Решение для наиболее общего случая, когда учитывается влияние всех трех механизмов на развитие трещины ГРП – коэффициента трещиностойкости, вязкости жидкости и фильтрационной утечки, получено в работе~\cite{Mitchell1}, в которой использован многопараметрический метод сингулярных возмущений. Для полученного решения рассмотрены предельные режимы малой и большой вязкости жидкости гидроразрыва. Во всех указанных работах используется приближение теории смазки~\cite{loic1}. Как следствие, локальные модели представляют собой системы уравнений параболического типа.



В данной работе рассматривается обобщение локальной модели PKN, учитывающее конвективные и нестационарные слагаемые в уравнении движения жидкости.

%%%%%%%%%%%%%%%%%%%%%%%%%%%%%%%%%%%%%%%%%%%%%%%%%%%%%%%%%%%%%%%%%%%%%%%%%%%%%%%%%%%%%%%%%%%%%%%%%%%%%%%%%%%%%%%%%%%%%%%%%%%%%%%%%%%%%%%%%%%%%%%%%%%%%%%%%%%%%%



\Section{Физическое обоснование модели}\label{sec:fiz_osn}

В представленном выше обзоре рассмотрены только работы, в которых жидкости разрыва являются ньютоновскими жидкостями. Это связано с тем, что во многих технологических процессах, связанных с ГРП, в качестве жидкости разрыва используется вода. К таким процессам относятся: 1) первая часть тестового (калибровочного) ГРП, в которой при скоростях закачки гидроразрыва пласта из скважины вытесняется пластовая вода линейным (не сшитым) гелем. В этом случае в течение нескольких минут вытесняемой водой создается трещина ГРП; 2) ГРП в низкопроницаемых нефтенасыщенных пластах; 3) явление авто - ГРП – самопроизвольное распространение магистральной трещины от нагнетательной скважины.



Во всех перечисленных процессах вследствие малой вязкости воды в трещине ГРП могут существовать достаточно протяженные области, в которых силы инерции превосходят или имеют один порядок с вязкими силами. Выполним оценку числа Рейнольдса. Расход жидкости при ГРП равен $2 \div 4$ $\text{м}^3/\text{мин}$. При высоте трещины $h = 10 \div 30\, \text{м}$ и ширине трещины у скважины (порядка диаметра перфорационных отверстий) $w = 0,01 \, \text{м}$ скорость течения в начале трещины равна $v =(0,11 \div 0,66)\, \text{м}/\text{с}$. Число Рейнольдса равно  $Re =(1100 \div 6600)$. Если конечная длина (полудлина, если трещина симметрична) трещины равна $L =(10 \div 50)\, \text{м}$, то соотношение инерционных и вязких сил для трещины в целом оценивается приведенным числом Рейнольдса $Re_L =Re \cdot w/L=(0,2 \div 6,6)$. Эта оценка показывает, что нельзя пренебрегать силами инерции в случае, когда жидкостью разрыва является вода.

	

Характерное время тестового ГРП (наименее продолжительного из процессов 1 - 3) составляет $T =180\, \text{с}$. Если характерное расстояние от стенок трещины до 

границ пласта порядка $L_0 = 200 \, \text{м}$, а скорость продольной волны в упругом пласте равна $a_0 = 4000\, \text{м}/\text{с}$, то характерное время пробега звуковой волны оценивается величиной $t_0 = L_0/a_0 = 0,05 \, \text{с}$. Отсюда следует оценка $T>> t_0$. Следовательно, процесс упругого деформирования пласта при закачке жидкости в трещину можно считать квазистационарным.



\Section{Квазиодномерная модель течения жидкости в вертикальной трещине}

Рассматривается процесс распространения вертикальной трещины постоянной высоты $h$ в насыщенном пористом пласте с коэффициентом проницаемости $k$ при нагнетании несжимаемой ньютоновской жидкости в пористый массив через вертикальную скважину. Предполагается, что трещина расположена симметрично относительно скважины.



Как и в модели PKN предполагается, что высота трещины много больше ее ширины $\delta$ и много меньше ее длины $L$, т.е. выполняется неравенство $\delta << h << L$. Кроме того, предполагается, что пласт является линейно-упругим, а по длине трещины условия напряженного состояния изменяются незначительно, что выполняется, если изменение давления жидкости внутри трещины много меньше характерного напряжения в пласте. Это дает возможность считать, что в плоскостях, перпендикулярных направлению трещины приближенно выполняется условие плоской деформации, а сечения не влияют друг на друга (гипотеза плоских сечений~\cite{Perkins_Kern, Nordgren1}).



Принимается, что ось $0X$ направлена вдоль трещины, ось $0Y$ направлена перпендикулярно оси $0X$ в направлении изменения ширины трещины $\delta$, а ось $0Z$ направлена по высоте трещины. Далее считается, что трещина расположена симметрично относительно плоскости $X0Y$, следовательно, верхний край трещины имеет координату $z = h/2$, а нижний край – координату $z = -h/2$.



В предыдущем пункте~\ref{sec:fiz_osn} показано, что процесс упругого деформирования можно считать квазистационарным. Следовательно, можно использовать известное стационарное решение для ширины $\delta$ симметричной трещины, находящейся в условиях плоской деформации в полубесконечном пространстве, на контур которой действует постоянное давление~\cite{Sneddon1}. При введенных выше допущениях для каждого поперечного сечения трещины ГРП можно приближенно написать:

\begin{equation}

\delta \left(x,z,t\right) = \frac{2\left(1-\nu\right)}{G} \left(p\left(x,t\right)-\sigma\right)\sqrt{\frac{h^2}{4}-z^2},\label{eq:2_1}

\end{equation}

где $\nu,\,G$ – соответственно коэффициент Пуассона и модуль сдвига породы; $p$ – давление в трещине; $\sigma$ – постоянное минимальное горизонтальное напряжение в пласте, в направлении которого трещина раскрывается при гидравлическом разрыве.



Интегрирование~(\ref{eq:2_1}) по высоте пласта позволяет найти площадь поперечного сечения как функцию координаты $x$ и времени $t$:

\begin{equation}

S\left(x,t\right) = \frac{\pi\left(1-\nu\right) h^2}{4G} \left(p\left(x,t\right)-\sigma\right).\label{eq:2_2}

\end{equation}

Если $S\left(x,t\right)$ - переменное поперечное сечение трещины, то с учетом несжимаемости жидкости можно ввести осредненные по сечению скорость, квадрат скорости и давление согласно выражениям

\begin{align}

\left\langle v\right\rangle &= \frac{1}{S} \int\limits_{S} vds,& \left(1+\beta\right)\left\langle v\right\rangle^2 &= \frac{1}{S} \int\limits_{S} v^2ds, & 

\left\langle p\right\rangle &= \frac{1}{S} \int\limits_{S} pds. \label{eq:2_3}

\end{align}

В подынтегральных выражениях в~(\ref{eq:2_3}) записаны истинные скорости и давления, а $\beta$ – поправка Кориолиса на неравномерное распределение истиной скорости по сечению~\cite{Charniy1}. Для стационарных ламинарных течений с параболическим профилем скорости $\beta = 1/3$, для квазистационарных турбулентных течений $\beta = 0$. При неустановившихся течениях поправка Кориолиса является функцией времени.



С учетом~(\ref{eq:2_3}) уравнения изменения массы и количества движения несжимаемой жидкости с постоянной плотностью $\rho$ в одномерном потоке в канале переменного сечения $S\left(x,t\right)$ с малой кривизной (необязательно вертикального) и проницаемыми стенками имеют вид:

\begin{align}

& \frac{\partial \left(\rho S\right)}{\partial t} + \frac{\partial \left(\rho \left\langle v\right\rangle  S\right)}{\partial x} = - \rho u_L \Pi, \label{eq:2_4} \\

& \frac{\partial \left(\rho \left\langle v\right\rangle  S\right)}{\partial t} + \frac{\partial \left(\left(1+\beta\right) \rho \left\langle v\right\rangle^2 S\right)}{\partial x} = 

-S \frac{\partial \left\langle p\right\rangle}{\partial x} - \Pi \tau - \rho g S \sin \alpha, \label{eq:2_5}

\end{align}

где $u_L$ нормальная стенкам канала скорость фильтрационной утечки; $\Pi$ – смоченный периметр; $\alpha$ - угол наклона канала к горизонту; $\tau$ - напряжение трения на стенке канала; $g$ ускорение силы тяжести. Как и в работе~\cite{Iv_Sm1} можно заменить произвольное сечение $S\left(x,t\right)$, прямоугольным сечением с площадью равной~(\ref{eq:2_2}):

\begin{equation}

S\left(x,t\right) = w\left(x,t\right)h .\label{eq:2_6}

\end{equation}



Из~(\ref{eq:2_2}) и~(\ref{eq:2_6}) следует локальная связь между давлением в трещине и ее раскрытием $w$ – осредненной шириной трещины $\delta$ по ее высоте $h$:

\begin{align}

p - \sigma &= bw, & b & = \frac{4G}{\pi h \left(1-\nu\right)} >0, \label{eq:2_7}

\end{align}

где $b$ – жесткость трещины.



С учетом~(\ref{eq:2_7}) система уравнений~(\ref{eq:2_4}) –~(\ref{eq:2_5}) запишется в дивергентном виде:

\begin{align}

& \frac{\partial \left(w\right)}{\partial t} + \frac{\partial \left(wv\right)}{\partial x} = - 2 u_L, \label{eq:2_8} \\

& \frac{\partial \left(wv \right)}{\partial t} + \frac{\partial \left(\left(1+\beta\right) wv^2 + bw^2/2\rho  \right)}{\partial x} = -2\frac{\tau}{\rho}, \label{eq:2_9}

\end{align}

где учтено, что трещина вертикальна, а смоченный периметр, в силу допущений модели, равен $\Pi\left(x,t\right) = 2 w\left(x,t\right) + 2h \approx 2h$. Скобки осреднения в уравнениях опущены.



Также систему$~(\ref{eq:2_8}) - ~(\ref{eq:2_9})$ можно представить в неконсервативном виде:

\begin{align}

A \frac{\partial \boldsymbol{U}}{\partial t} + B \frac{\partial \boldsymbol{U}}{\partial x} = \boldsymbol{d}, $  $ \boldsymbol{U}  = \left(w,v\right)^T,$  $ A &= \left(\begin{matrix}1 & 0 \\ v & w \end{matrix}\right), \label{eq:2_10} \\

\left( \begin{matrix}v & w \\ \left[\left(1+\beta\right) v^2 + \frac{bw}{\rho}\right] & \left[2\left(1+\beta\right)wv\right] \end{matrix}\right) &= B,   \boldsymbol{d} = \left(-2u_L, -2 \frac{\tau}{\rho}\right)^T.\label{eq:2_11}

\end{align}



Отсюда следует, что данная система уравнений является квазилинейной. 



Для замыкания системы уравнений~(\ref{eq:2_8}) –~(\ref{eq:2_9}) или~(\ref{eq:2_10}) –~(\ref{eq:2_11}) необходимо задать скорость фильтрации жидкости через ее стенки и напряжение трения на стенках трещины.



Далее предполагается, что: 1) пласт бесконечен и в начальный момент времени пластовое (поровое) давление постоянно и равно $p_k$; 2) течение в пласте квазистационарно и подчиняется уравнениям упругого режима фильтрации; 3) в каждом сечении трещины фильтрационная утечка жидкости через стенки трещины не зависит от утечек в соседних сечениях и представляет собой прямолинейно-параллельный поток в пласте. При допущениях 1) – 3) применение автомодельного решения для прямолинейного течения в бесконечном линейно-упругом пласте, при заданном постоянном давлении на границе полупространства (течение от нагнетательной галереи)~\cite{Smirn_Tag1}, приводит к двухмерному распределению давления $p\left(x,y,t\right)$ в пласте:

\begin{equation}

p\left(x,y,t\right) = p \left(x,t\right) + \left(p_k - p \left(x,t\right)\right) erf \left(\frac{y}{2 \sqrt{\chi t}}\right) , \label{eq:2_12}

\end{equation}

где $\chi = k/\mu \beta_*$ - коэффициент пъезопроводности породы; $\beta_* = \beta_l m + \beta_s$ – коэффициент упругой сжимаемости пласта; $\beta_l m,\, \beta_s$ – соответственно коэффициенты упругой сжимаемости пластовой жидкости и скелета породы.



С учетом~(\ref{eq:2_12}) скорость фильтрационной утечки жидкости через стенки трещины принимается в виде:

\begin{equation}

u_L = \left.-\frac{k}{\mu} \frac{\partial p\left(x,y,t\right)}{\partial y}\right|_{y=0} = \frac{k}{\mu} \frac{\left(p \left(x,t\right) - p_k\right)}{\sqrt{\pi \chi \left(t-t_L\right)}}, \label{eq:2_13}

\end{equation}

где $k$ – проницаемость пласта; $\mu$  вязкость закачиваемой жидкости. В формуле~(\ref{eq:2_13}) учтено запаздывание утечки вследствие распространения фронта трещины. В~(\ref{eq:2_13}) , $t_L = f^{-1} (L)$ а $L= f(t)$  – закон движения фронта трещины.



Нужно отметить, что допущения 1) – 3) являются модельными и вместо~(\ref{eq:2_13}) можно использовать другие эмпирические или аналитические выражения для скорости фильтрационной утечки.

При известной поправке Кориолиса ?, полученная система уравнений~(\ref{eq:2_8}) –~(\ref{eq:2_9}) замыкается заданием напряжения трения $\tau$ и коэффициента трения $C_w$ при ламинарном режиме ($Re \leq 2000$) течения~\cite{Idel4ik1}:

\begin{align}

&\tau = C_w \frac{\rho v^2 }{2}, C_w  = \frac{\lambda}{4} = \frac{k_0 \cdot 64}{4 Re} = \frac{24}{Re}, \label{eq:2_14}

\end{align}

где $k_0 = 3/2$ – поправочный коэффициент для каналов прямоугольного сечения, когда одна сторона прямоугольника много больше другой его стороны; $\lambda$ - коэффициент гидравлического сопротивления. В уравнении~(\ref{eq:2_14}) число Рейнольдса $Re$, определяется через гидравлический диаметр $D_H$ сечения трещины:

\begin{align}

Re &= \frac{\rho v D_H }{\mu} , & D_H & = \frac{4S}{\Pi} = \frac{4wh }{2 \left(w+h\right)} \approx 2w. \label{eq:2_15}

\end{align}



При турбулентном режиме~($Re>2000$) течения~\cite{Idel4ik1}:

\begin{align}

&\tau = C_w \frac{\rho v^2 }{2} ,  C_w  = \frac{\lambda_1}{4}= \frac{\lambda k_0}{4} = \frac{k_0 \cdot 0.3164}{4 Re^{0.25}} = \frac{0.08701}{Re^{0.25}}, \label{eq:2_16}

\end{align}

где пренебрегается шероховатостью стенок трещины. $\lambda_1$ коэффициент гидравлического сопротивления при турбулентном режиме течения в канале прямоугольного сечения; $\lambda$ - коэффициент гидравлического сопротивления при турбулентном режиме течения в круглой трубе (формула Блазиуса); $k_0 = 1,1$ – поправочный коэффициент для турбулентных потоков в каналах прямоугольного сечения, когда одна сторона прямоугольника много больше другой его стороны.



\Section{Характеристики, соотношения на характеристиках и инварианты Римана} \label{sec:3}

Матрица $A$ в~(\ref{eq:2_10}) не вырождена при условии $w \neq 0$. Умножение векторного уравнения~(\ref{eq:2_10}) слева на обратную матрицу $A^{-1}$, дает новое представление данного уравнения

\begin{align}

&\frac{\partial \boldsymbol{\boldsymbol{U}}}{\partial t} + C \frac{\partial \boldsymbol{\boldsymbol{U}}}{\partial x} = \boldsymbol{f}, C = A^{-1} B = \left( \begin{matrix}v & w \\ \left(\beta v^2 + \rho^{-1} bw)/w\right  & \left\left(2\beta+1\right)v\right \end{matrix}\right), \label{eq:3_1} \\

&\boldsymbol{\boldsymbol{f}}  = A^{-1} \boldsymbol{d} = \left(-2u_L, 2 \frac{\left(u_L v - \tau / \rho\right)}{w}\right)^T, V{\boldsymbol{U}} = \left(w,v\right)^T. \label{eq:3_2}

\end{align}



Согласно общей теории квазилинейных уравнений~\cite{kulik1,rozhd1,ovsyan1,Chern1}, вычисляются собственные значения матрицы $C$, что при условии $w \neq 0$ дает два вещественных собственных значения на решении $\boldsymbol{U}$:

\begin{equation}

\lambda_{\pm} = \left(\beta + 1\right)v \pm \sqrt{\left(\beta^2 + \beta\right) v^2 + \frac{bw}{\rho}}.

\label{eq:3_3}

\end{equation}



Дифференциальные уравнения характеристических кривых имеют вид:

\begin{align}

C_+:& & \frac{dx}{dt} & = \left(\beta + 1\right)v + \sqrt{\left(\beta^2 + \beta\right) v^2 + \frac{bw}{\rho}}, \label{eq:3_4} \\

C_-:& & \frac{dx}{dt} & = \left(\beta + 1\right)v - \sqrt{\left(\beta^2 + \beta\right) v^2 + \frac{bw}{\rho}}. \label{eq:3_5}

\end{align}



Находятся левые собственные векторы:

\begin{align}

\boldsymbol{l_1} & = \left(\frac{1}{w}\left(-\beta v + \sqrt{\left(\beta^2 + \beta\right) v^2 + \frac{bw}{\rho}}\right),\, 1 \right)^T , \label{eq:3_6} \\

\boldsymbol{l_2} & = \left(\frac{1}{w}\left(-\beta v - \sqrt{\left(\beta^2 + \beta\right) v^2 + \frac{bw}{\rho}}\right),\, 1 \right)^T . \label{eq:3_7} 

\end{align}



Матрица $\Omega_L$, составленная из левых собственных векторов не вырождена при условии $w \neq 0$, поскольку

\begin{equation}

\det \Omega_L = \det \left(\begin{matrix}l_{11} & l_{12} \\ l_{21} & l_{22} \end{matrix}\right) = \frac{2}{w}\sqrt{\left(\beta^2 + \beta\right) v^2 + \frac{bw}{\rho}} \neq 0,

\label{eq:3_8}

\end{equation}

где номер строки обозначает номер собственного вектора, а номер столбца – его координату. Таким образом, система уравнений~(\ref{eq:3_1}),~(\ref{eq:3_2}) является строго гиперболической системой~\cite{kulik1}.



Соотношения на характеристиках имеют вид:

\begin{align}

&\frac{1}{w}\left(-\beta v + \sqrt{\left(\beta^2 + \beta\right) v^2 + \frac{bw}{\rho}}\right) D_+ w + D_+ v  = \nonumber \\

& = -\frac{2}{w}\left(-\beta v + \sqrt{\left(\beta^2 + \beta\right) v^2 + \frac{bw}{\rho}}\right) u_L + \frac{2}{w} \left(u_L v - \frac{\tau}{\rho}\right), \label{eq:3_9} \\

&\frac{1}{w}\left(-\beta v - \sqrt{\left(\beta^2 + \beta\right) v^2 + \frac{bw}{\rho}}\right) D_- w + D_- v  = \nonumber \\

& = \frac{2}{w}\left(\beta v + \sqrt{\left(\beta^2 + \beta\right) v^2 + \frac{bw}{\rho}}\right) u_L + \frac{2}{w} \left(u_L v - \frac{\tau}{\rho}\right), \label{eq:3_10} \\

\end{align}

где введены операторы внутреннего дифференцирования вдоль характеристик~\cite{ovsyan1}:

\begin{align}

 D_+&= \partial_t + \lambda_+\partial_x, & D_-&= \partial_t + \lambda_-\partial_x, \nonumber

\end{align}

а в правых частях~(\ref{eq:3_9}) и~(\ref{eq:3_10}) скорость фильтрации в пласт и пристеночное трение выражаются через искомые функции по формулам~(\ref{eq:2_13})-(\ref{eq:2_15}).



Как известно~\cite{rozhd1}, квазилинейную гиперболическую систему из двух уравнений можно записать в инвариантах Римана. Однако, при любом $\beta >0$ инварианты Римана системы~(\ref{eq:3_9}),~(\ref{eq:3_10}) не выражаются в элементарных функциях. Это удается сделать, если пренебречь поправкой Кориолиса, т.е. положить $\beta = 0$. В этом случае соотношения вдоль характеристик запишутся в виде:

\begin{align}

 D_+I_1&= -\frac{2}{w} \sqrt{\frac{bw}{\rho}} u_L + \frac{2}{w} \left(u_L v - \frac{\tau}{\rho}\right), \label{eq:3_11} \\

 D_-I_2&= \frac{2}{w} \sqrt{\frac{bw}{\rho}} u_L + \frac{2}{w} \left(u_L v - \frac{\tau}{\rho}\right), \label{eq:3_12}

\end{align}

где инварианты Римана суть

\begin{align}

 I_{1,2}&= v \pm 2 \sqrt{\frac{bw}{\rho}}. \label{eq:3_13}

\end{align}



Правые части соотношений~(\ref{eq:3_11}) и~(\ref{eq:3_12}) выражаются через инварианты Римана с учетом~(\ref{eq:2_7}),~(\ref{eq:2_13})-(\ref{eq:2_16}) и следующих из~(\ref{eq:3_13}) представлений искомых функций через эти инварианты

\begin{align}

w & = \rho \left(I_1 - I_2\right)^2 \frac{1}{16b}, & v& = \frac{\left(I_1+I_2\right)}{2}. \nonumber

\end{align}



Дифференциальные уравнения характеристик~(\ref{eq:3_4}),~(\ref{eq:3_5}) при $\beta = 0$ имеют вид:

\begin{align}

C_+:& & \frac{dx}{dt} & = v + \sqrt{\frac{bw}{\rho}}, & C_-:& & \frac{dx}{dt} & = v - \sqrt{\frac{bw}{\rho}}, \label{eq:3_14}

\end{align}

и величину

\begin{align}

&a= \sqrt{\frac{bw}{\rho}} = \sqrt{\frac{\left(p-\sigma\right)}{\rho}} \label{eq:3_15}

\end{align}

можно трактовать как скорость малых возмущений раскрытия трещины или давления (играет роль скорости звука). Отметим, что~(\ref{eq:3_15}) совпадает со скоростью малых возмущений давления в трубке с упругими стенками, вычисленной для линейной связи~(\ref{eq:2_7}) через растяжимость стенок трубки~\cite{pedli1}.



Для наглядности на рисунке~\ref{r1} показаны зависимости местной скорости $a$ малых возмущений давления~(\ref{eq:3_15}) несжимаемой жидкости с плотностью $\rho = 1000 \text{кг}/\text{м}^3$ от раскрытия $w$ трещины для различных значений высоты трещины $h$. Модуль сдвига равен $G = 25 \text{ГПа}$, коэффициент Пуассона породы равен $\nu = 0,13$ (песчаник).

\begin{figure} [h!] 

  \begin{center}

	\includegraphics{pics/r1.png}	

	\end{center}

  \caption{Зависимость местной скорости малых возмущений давления $a\, (\text{м/с})$ от раскрытия трещины $w\, (\text{м})$ для различных значений высоты трещины $h\, (\text{м})$. 1– 10; 2–20; 3–30; 4–40; 5–50.}  

  \label{r1}  

\end{figure}

Очевидно, что уравнения модели развития трещины гидроразрыва и все предыдущие соотношения можно легко переписать в терминах скорости и давления жидкости.



\Section{Граничные условия для трещины ГРП. Эволюционность границ трещины. Постановка начально-краевой задачи}

Далее рассматриваются вопросы о граничных и начальных условиях уравнений представленной модели развития трещины гидроразрыва пласта. Вопросы корректности начально-краевых задач для систем квазилинейных гиперболических уравнений и уравнений газовой динамики рассмотрены в~\cite{rozhd1}. Вопросы о количестве граничных условий на подвижных и неподвижных границах для уравнений газовой динамики рассмотрены в~\cite{rozhd1,Zver_smirn1,Chern1}, а для квазилинейных уравнений в~\cite{kulik1,rozhd1,Zver_smirn1}.



Трещина ГРП образуется в области перфорационных отверстий вертикальной скважины, т.е. входная граница является неподвижной. Обычно на практике при инициации трещины ГРП задается постоянный расход $Q_0$ закачиваемой жидкости, поэтому граничным условием на входе трещины ГРП является нелинейное граничное условие:

\begin{equation}

Q_0 - hwv = 0.

\label{eq:4_1}

\end{equation}



Рассмотрим вопрос об эволюционности входной границы трещины ГРП. Пусть в области $x > 0$ в начальный момент времени заданы начальные условия. Из области течения $x > 0$ к неподвижной границе $x = 0$ приходит одна характеристика с характеристической скоростью $\lambda_- = v-a<0$. От входной границы в область течения уходит характеристика второго семейства $\lambda_+ = v+a>0$. Входная граница будет эволюционной, если на ней задано одно граничное условие~(\ref{eq:4_1}), так как в этом случае на границе определяются раскрытие $w$ и скорость течения $v$. 



Предположим, что фронт трещины представляет собой сильный разрыв ее раскрытия (давления). 

В случае пренебрежения неоднородностью профиля скорости $\left(\beta = 0\right)$ применение интегральных законов сохранения массы и количества движения~\cite{Chern1} к установившемуся квазиодномерному течению в канале с упругими стенками с учетом сил давления, возникающих на скачке в системе координат, связанной с разрывом, приводит к следующим условиям на скачке:     

\begin{align}

S_2 v_2 &= S_1 v_1 -\int\limits_{1}^{2} \Pi u_L dx, \nonumber \\

S_2 \left(p+\rho v_2^2\right) & = S_1 \left(p+\rho v_1^2\right) + \int\limits_{1}^{2} p(S) dS -\int\limits_{1}^{2}\Pi \tau dx - \rho \int\limits_{1}^{2}\Pi v u_L dx,\nonumber

\end{align}

где считается, что скачок раскрытия имеет некоторую ширину. Индексом 1 обозначены параметры перед разрывом, а индексом 2 – параметры за разрывом. Для простоты предполагается, что плотности закачиваемой и пластовой жидкостей совпадают.



При пренебрежении фильтрационной утечкой и вязкостью жидкости эти соотношения совпадают с условиями на ударной волне в идеальной несжимаемой жидкости~\cite{pedli1} в трубке с упругими стенками, если связь $p = p(S)$ имеет место на ударной волне. 



Поскольку на практике выполняется неравенство $u_L << v$, то последними интегралами в правых частях законов сохранения можно пренебречь. Если, кроме того, второй интеграл в правой части закона сохранения количества движения много меньше первого, то в системе координат наблюдателя в случае линейной связи~(\ref{eq:2_7}) условия на скачке раскрытия примут вид:

\begin{align}

w_1 \left(v_1 - D\right) &=w_2 \left(v_2 - D\right) = m , \label{eq:4_2} \\

\rho w_1 \left(v_1 - D\right)^2 + \frac{b}{2} w_1^2 &=\rho w_2 \left(v_2 - D\right)^2 + \frac{b}{2} w_2^2 . \label{eq:4_3} 

\end{align}



В~(\ref{eq:4_2}),~(\ref{eq:4_3}) $D$ – скорость скачка; $m$ – поток массы нагнетаемой жидкости через разрыв. 



Далее рассматривается вопрос об эволюционности фронта трещины ГРП, который по предположению является сильным разрывом раскрытия трещины. В этом случае $m \neq 0$ и жидкость течет через разрыв. Рассмотрим течение жидкости разрыва в области $x \geq 0$  с подвижной внутренней границей $x = x(t)$ – сильным скачком раскрытия трещины ГРП, если заданы начальные условия и краевое условие~(\ref{eq:4_1}). На скачке заданы пять неизвестных величин $w_1,\, v_1,\, w_2,\, v_2,\, D$. Для того, чтобы разрыв раскрытия трещины был эволюционным, число заданных граничных условий на нем должно равняться $N= 5-s$, где $s$ – общее число приходящих на разрыв характеристик из областей течения слева и справа от разрыва. 



В~\cite{pedli1} показано, что в ударных волнах в упругой трубке аналогично ударным волнам в газовой динамике, скорость течения жидкости относительно скачка перед скачком сверхкритическая $\left(\left|v_1-D\right|> a_1\right)$, а за скачком докритическая $\left(\left|v_2-D\right|< a_2\right)$. Это означает, что из области слева от разрыва на скачок приходит одна характеристика со скоростью  $\lambda_+ = v_2 - D + a>0$, а из области справа от разрыва на скачок приходят две характеристики со скоростями $\lambda_+ = v_1 - D + a<0$ и $\lambda_- = v_1 - D - a<0$. Итак, $s=3,\, N=2,$ и задание двух законов сохранения~(\ref{eq:4_2}) и~(\ref{eq:4_3}) на скачке достаточно для его эволюционности. 



В случае $m = 0$ жидкость не течет через разрыв. Из~(\ref{eq:4_2}) и~(\ref{eq:4_3}) следует, что в этом случае течение на фронте трещины ГРП является непрерывным, поскольку выполняются равенства $w_1 =w_2, \, v_1 = v_2$. Скорость скачка в этом случае является скоростью распространения фронта трещины.



Поскольку в данной модели развитие трещины разрыва породы не моделируется, то предлагается следующая идеализированная постановка задачи о развитии трещины ГРП: при $t>0$ в области $x \geq 0$  найти обобщенное решение~\cite{rozhd1} гиперболической системы уравнений~(\ref{eq:3_1}),~(\ref{eq:3_2}) с граничным условием~(\ref{eq:4_1}) на входной границе, условиями на бесконечности

\begin{align}

\lim\limits_{x\to \infty} w\left(x,t\right) &= w_0 \neq 0, &\lim\limits_{x\to \infty} v\left(x,t\right) &= 0  \nonumber

\end{align}

и начальными условиями

\begin{align}

w\left(x,0\right) &= w_0 \neq 0, & v\left(x,0\right) &= 0 . \nonumber

\end{align}



Требование неравенства нулю раскрытия трещины в начальный момент времени следует из условия~(\ref{eq:3_8}) гиперболичности системы уравнений~(\ref{eq:3_1}),~(\ref{eq:3_2}).



\Section{Аналогия с газовой динамикой}

Для численного решения полученной системы уравнений методами типа Годунова~\cite{kulik1} для нее необходимо решить задачу о распаде произвольного начального разрыва в пренебрежении диссипативными слагаемыми – трением в трещине ГРП и фильтрационной утечкой. В этом случае при $\beta = 0$ система уравнений~(\ref{eq:3_1}),~(\ref{eq:3_2}) с учетом~(\ref{eq:3_15}) примет вид:

\begin{align}

\frac{\partial w}{\partial t} + v\frac{\partial w}{\partial x}+w\frac{\partial v}{\partial x} =0, \label{eq:5_1} \\

\frac{\partial v}{\partial t} + \frac{a^2}{w}\frac{\partial w}{\partial x}+v\frac{\partial v}{\partial x} =0. \label{eq:5_2}

\end{align}



Уравнения~(\ref{eq:5_1}) и~(\ref{eq:5_2}) аналогичны уравнениям газовой динамики в гидравлическом приближении~\cite{Zver_smirn1}, если в последних уравнениях газ считать несжимаемой жидкостью. Можно также заметить, что существует аналогия с уравнениями одномерной газовой динамики с изоэнтропическими плоскими волнами~\cite{rozhd1,ovsyan1,Chern1}, если учесть, что в уравнениях~(\ref{eq:5_1}) и~(\ref{eq:5_2}) раскрытие трещины играет одновременно как роль плотности, так и роль давления. 

Применительно к системе уравнений~(\ref{eq:5_1}),~(\ref{eq:5_2}) можно ввести определения в терминах раскрытия трещины $w$ и скорости жидкости $v$, инварианты Римана $I_1,\, I_2$ ~\cite{pedli1}, а также получить решения типа простых волн (ПВ): волн Римана (ВР) и центрированных волн и доказать утверждения аналогичные утверждениям в одномерной газовой динамике с изоэнтропическими плоскими 

волнами~\cite{ovsyan1}.

\begin{enumerate}

	\item Вырожденная ПВ системы~(\ref{eq:5_1}),~(\ref{eq:5_2}) есть постоянное течение. Ее линиями уровня являются траектории частиц жидкости, покрывающие плоскость событий. Линиями уровня невырожденной ПВ являются прямолинейные характеристики~(\ref{eq:3_14}), покрывающие плоскость событий.

	\item Течение, удовлетворяющее системы уравнений~(\ref{eq:5_1}),~(\ref{eq:5_2}) является невырожденной ПВ тогда и только тогда, когда один из инвариантов Римана тождественно постоянен. Если $I_1 \equiv const$, то линии уровня ПВ – прямолинейные характеристики $C_-$. Если $I_2 \equiv const$, то линиями уровня ПВ являются прямолинейные характеристики $C_+$.

	\item (Достаточный признак ПВ) Если в непрерывном течении имеется характеристика $C_+\left(C_-\right)$, вдоль которой постоянны $w$ и $v$, причем $w \neq 0$, то с каждой ее стороны течение является либо постоянным, либо ВР.

	\item Невырожденная ВР является волной сжатия (разрежения) тогда и только тогда, когда ее прямолинейные характеристики с ростом времени сходятся (расходятся) на плоскости событий. В волне сжатия частицы разгоняются, а в волне разрежения – тормозятся.

\end{enumerate}



\Section[n]{Заключение}

В работе предложена математическая модель развития трещины ГРП гиперболического типа, являющаяся обобщением модели PKN. Исследована эволюционность границ трещины ГРП. Предложена постановка начально-краевой задачи для представленных уравнений модели. Показано, что при пренебрежении диссипативными слагаемыми в представленной модели можно построить теорию простых волн, аналогичную теории одномерной газовой динамики с изоэнтропическими плоскими волнами.



{\thanks{Работа выполнена при финансовой поддержке Минобрнауки РФ в рамках базовой части государственного задания организациям высшего образования в 2016 г. и проекта РФФИ $($ \rm \No~$14{-}01{-}97012\,p).$ \\ Булгакова Гузель Талгатовна: http://orcid.org/0000-0001-8030-1791\\

Ильясов Айдар Мартисович:   http://orcid.org/0000-0002-4399-9040

}}



%% Благодарности, гранты и прочее

\newpage

\begin{thebibliography}{9}



%%%%%%%%%%%%%%%%%%%%%%%%%%%%%%%%%%%%%%%%%%%%

\RBibitem{Zheltov-Christ}

\paper О гидравлическом разрыве нефтеносного пласта

\by   Желтов~Ю.\,Н., Христианович~С.\,А.

\jour Изв. АН СССР. ОТН.

\yr 1955

\issue 6

\pages 3--41



\RBibitem{Mush}

\by Мусхелишвили~Н.\,Н.

\book Некоторые oсновные задaчи математической теории упругости

\publaddr М.

\publ Наука

\yr 1966

\finalinfo 707~с



\Bibitem{Perkins_Kern}

\paper Width of hydraulic fractures 

\by   Perkins~T.\,K., Kern.~L.\,R.

\jour Journal of Petroleum Technology

\yr 1961

\issue 13, 4

\pages 937--949



\Bibitem{Nordgren1}

\paper Propogation of a vertical hydraulic fracture

\by   Nordgren~R.\,P.

\jour Society of Petroleum Engineers Journal

\yr 1972

\issue 12, 4

\pages 306--314



\RBibitem{Iv_Sm1}

\paper Формирование трещины гидроразрыва в пористой среде

\by   Ивашнев~О.\,Е.,~Смирнов~Н.\,Н.

\jour Вестн. МГУ. Математика, механика

\yr 2003

\issue 6

\pages 28--36



\RBibitem{Gord_Z1}

\paper Автомодельное решение задачи о глубокопроникающем гидравлическом разрыве пласта

\by   Гордеев~Ю.\,Н.,~Зазовский~А.\,Ф.

\jour Известия РАН. МТТ

\yr 1991

\issue 5

\pages 119--131



\RBibitem{Smirn_Tag1}

\paper Автомодельные решения задачи о формировании трещины гидроразрыва в пористой среде

\by   Смирнов~Н.\,Н.,~Тагирова~В.\,Р.

\jour Изв. РАН. МЖГ

\yr 2007

\issue 1

\pages 70--82





\RBibitem{Basn1}

\by Басниев~К.\,С.,~Кочина~И.\,Н.,~Максимов~В.\,М.

\book Подземная гидромеханика

\publaddr М.

\publ Недра

\yr 1993

\finalinfo 416~с



\RBibitem{Teod_tr1}

\paper Форма плоской трещины гидрoразрыва в упругой непроницаемой среде при различных скоростях закачки

\by   Теодорович~Э.\,В.,~Трофимов~А.\,А.,~Шумилин~И.\,Д.

\jour Изв. РАН. МЖГ

\yr 2011

\issue 4

\pages 252--261



\Bibitem{Garagash1}

\paper Plane-strain propagation of a hydraulic fracture: small toughness solution

\by   Garagash~D.\,I.,~Detournay~E.

\jour J. Appl. Mech.

\yr 2005

\issue 72

\pages 916-–928



\Bibitem{Garagash2}

\paper Plane strain propagation of a hydraulic fracture during injection and shut-in: large toughness solutions

\by   Garagash~D.\,I.

\jour Eng. Frac. Mech.

\yr 2006

\issue 73

\pages 456-–481



\Bibitem{Adachi1}

\paper Plane-strain propagation of a hydraulic fracture in a permeable rock

\by   Adachi~J.\,I.,~Detournay~E.

\jour Eng. Frac. Mech.

\yr 2008

\issue 75

\pages 4666-–4694



\Bibitem{Mitchell1}

\paper An asymptotic framework for finite hydraulic fractures including leak-off

\by   Mitchell~S.\,L,~Kuske~R,~Peirce~A.\,P.

\jour SIAM J. Appl. Math.

\yr 2007

\issue 67

\pages 346-–386



\RBibitem{loic1}

\by Лойцянский~Л.\,Г.

\book Механика жидкости и газа

\publaddr М.

\publ Дрофа

\yr 2003

\finalinfo 840~с



\RBibitem{Sneddon1}

\by Снеддон~И.\,Н.,~Берри~Д.\,С.

\book Классическая теория упругости

\publaddr М.

\publ Физматгиз

\yr 1961

\finalinfo 219~с





\RBibitem{Charniy1}

\by Чарный~И.\,А.

\book Неустановившееся движение реальной жидкости в трубах

\publaddr М.

\publ Недра

\yr 1975

\finalinfo 296~с





\RBibitem{Idel4ik1}

\by Идельчик~И.\,Е.

\book Справочник по гидравлическим сопротивлениям

\publaddr М.

\publ Машиностроение

\yr 1992

\finalinfo 672~с





\RBibitem{kulik1}

\by Куликовский~А.\,Г.,~Погорелов~Н.\,В.,~Семенов~А.\,Ю.

\book Математические вопросы численного решения гиперболических систем уравнений

\publaddr М.

\publ Физматлит

\yr 2001

\finalinfo 608~с



\RBibitem{rozhd1}

\by Рождественский~Б.\,Л.,~Яненко~Н.\,Н.

\book Системы квазилинейных уравнений и их приложения к газовой динамике

\publaddr М.

\publ Наука

\yr 1978

\finalinfo 688~с





\RBibitem{ovsyan1}

\by Овсянников~Л.\,В.

\book Лекции по основам газовой динамики

\publaddr М.-Иж-ск

\publ Инс-т. комп. иссл.

\yr 2003

\finalinfo 336~с





\RBibitem{pedli1}

\by Педли~Т.

\book Гидродинамика крупных кровеносных сосудов

\publaddr М.

\publ Мир

\yr 1983

\finalinfo 400~с



\RBibitem{Zver_smirn1}

\by Зверев~И.\,Н.,~Смирнов~Н.\,Н.

\book Газодинамика горения

\publaddr М.

\publ Изд-во МГУ

\yr 1987

\finalinfo 307~с





\RBibitem{Chern1}

\by Черный~Г.\,Г.

\book Газовая динамика

\publaddr М.

\publ Наука

\yr 1988

\finalinfo 424~с





\end{thebibliography}







\newpage







\chead[\fancyplain{}{\scriptsize \sl I\,l\,y\,a\,s\,o\,v~A.\,M.,

B\,u\,l\,g\,a\,k\,o\,v\,a~G.\,T.}]{\fancyplain{}{\scriptsize \sl  The quasi-one-dimensional hyperbolic model of hydraulic fracturing \dots}}



\makeengtitlem

\newpage

\begin{thebibliography}{10}



%%%%%%%%%%%%%%%%%%%%%%%%%%%%%%%%%%%%%%%%%%%%

\RBibitem{Zheltov-Christ}

\paper Formation of vertical fractures by means of highly viscous liquid

\by   Zheltov~Y.\,P. , Khristianovich~S.\,A.

\jour Proceedings of the Fourth World Petroleum Congress

\yr 1955

\issue 6

\pages 3--41



\RBibitem{Mush}

\by 2.	Muskhelishvili~N.\,N. 

\book Nekotorye osnovnye zadachi matematicheskoi teorii uprugosti

\publaddr М.

\publ Nauka

\yr 1966

\finalinfo 707~с



\Bibitem{Perkins_Kern}

\paper Width of hydraulic fractures 

\by   Perkins~T.\,K., Kern.~L.\,R.

\jour Journal of Petroleum Technology

\yr 1961

\issue 13, 4

\pages 937--949



\Bibitem{Nordgren1}

\paper Propogation of a vertical hydraulic fracture

\by   Nordgren~R.\,P.

\jour Society of Petroleum Engineers Journal

\yr 1972

\issue 12, 4

\pages 306--314



\RBibitem{Iv_Sm1}

\paper Formirovanie treshchiny gidrorazryva v poristoi srede 

\by   	Ivashnev~O.\,E., Smirnov~N.\,N

\jour Vestnik Moskovskogo universiteta. Matematika, mekhanika

\yr 2003

\issue 6

\pages 28--36



\RBibitem{Gord_Z1}

\paper Avtomodel'noe reshenie zadachi o glubokopronikaiushchem gidravlicheskom razryve plasta 

\by   	Gordeev~Ju.\, N., Zazovskey~A.\,F.

\jour Izvestia RAN. Mekhanika tverdogo tela. – Mechanics of Solids.

\yr 1991

\issue 5

\pages 119--131







\RBibitem{Smirn_Tag1}

\paper Self-similar solutions of the problem of formation of a hydraulic fracture in a porous medium 

\by   Smirnov~N.\,N., Tagirova~V.\,P.

\jour Fluid Dyn

\yr 2007

\issue 1

\pages 70--82





\RBibitem{Basn1}

\by Basniev~K.\,S., Kochina~I.\,N., Maksimov~V.\,M. 

\book Podzemnaia gidromekhanika

\publaddr Moscow

\publ Nedra

\yr 1993

\finalinfo 416~p



\RBibitem{Teod_tr1}

\paper Shape of a plane hydraulic fracture crack in an elastic impermeable medium at various injection rates 

\by   	Teodorovich~E.\,V., Trofimov~A.\,A. , Shumilin~I.\,D.  

\jour Fluid Dyn. 

\yr 2011

\issue 4

\pages 252--261



\Bibitem{Garagash1}

\paper Plane-strain propagation of a hydraulic fracture: small toughness solution

\by   Garagash~D.\,I.,~Detournay~E.

\jour J. Appl. Mech.

\yr 2005

\issue 72

\pages 916-–928



\Bibitem{Garagash2}

\paper Plane strain propagation of a hydraulic fracture during injection and shut-in: large toughness solutions

\by   Garagash~D.\,I.

\jour Eng. Frac. Mech.

\yr 2006

\issue 73

\pages 456-–481



\Bibitem{Adachi1}

\paper Plane-strain propagation of a hydraulic fracture in a permeable rock

\by   Adachi~J.\,I.,~Detournay~E.

\jour Eng. Frac. Mech.

\yr 2008

\issue 75

\pages 4666-–4694



\Bibitem{Mitchell1}

\paper An asymptotic framework for finite hydraulic fractures including leak-off

\by   Mitchell~S.\,L,~Kuske~R,~Peirce~A.\,P.

\jour SIAM J. Appl. Math.

\yr 2007

\issue 67

\pages 346-–386



\RBibitem{loic1}

\by 	Lojcansky~L.\,G.

\book Mekhanika zhidkosti i gaza

\publaddr Moscow

\publ Drofa

\yr 2003

\finalinfo 840~p



\RBibitem{Sneddon1}

\by 	Sneddon~J.\,N., Berry~D.\,S. 

\book The Classical Theory of Elasticity.

\publaddr Moscow

\publ  Springer

\yr 1961

\finalinfo 219~p





\RBibitem{Charniy1}

\by 	Charnyi~I.\,A. 

\book Neustanovivsheesia dvizhenie real'noi zhidkosti v trubakh

\publaddr Moscow

\publ Nedra

\yr 1975

\finalinfo 296~p





\RBibitem{Idel4ik1}

\by 	Idelchik~I.\,E.

\book Handbook of Hydraulic Resistance

\publaddr New York

\yr 2008







\RBibitem{kulik1}

\by 	Kulikovskey~A.\,G., Pogorelov~N.\,B., Semenov~A.\,Ju. 

\book Matematicheskie voprosy chislennogo resheniia giperbolicheskikh sistem uravnenii

\publaddr Moscow

\publ Fizmalit

\yr 2001

\finalinfo 608~p



\RBibitem{rozhd1}

\by 	Rogdestvenskii~B.\,L., Yanenko~N.\,N. 

\book Sistemy kvazilineinykh uravnenii i ikh prilozheniia k gazovoi dinamike

\publaddr Moscow

\publ Nauka

\yr 1978

\finalinfo 688~p





\RBibitem{ovsyan1}

\by 	Ovsiannikov~L.\,B.

\book Lektsii po osnovam gazovoi dinamiki

\publaddr Moscow – Izhevsk

\publ Institut komp'iuternykh issledovanii.

\yr 2003

\finalinfo 336~p





\RBibitem{pedli1}

\by 	Pedley~T.\,J.

\book The fluid mechanics of large blood vessels

\publaddr М.

\publ Cambridge University Press

\yr 1980

\finalinfo 400~с



\RBibitem{Zver_smirn1}

\by 	Zverev~I.\,N., Cmirnov~N.\,N.

\book Gazidinamika gorenia

\publaddr Moscow

\publ Mos. St. University

\yr 1987

\finalinfo 307~p





\RBibitem{Chern1}

\by 	Chernyi~G.\,G. 

\book Gazovaia dinamika

\publaddr Moscow

\publ Nauka

\yr 1988

\finalinfo 424~p





\end{thebibliography}



\date{12/XI/2016} % Дата поступления статьи в редакцию.

\revisiondate{12/XI/2016} % В окончательном виде



% \newpage

%

% \tableofengcontents

%

\end{document}

